% Options for packages loaded elsewhere
\PassOptionsToPackage{unicode}{hyperref}
\PassOptionsToPackage{hyphens}{url}
\documentclass[
]{article}
\usepackage{xcolor}
\usepackage{amsmath,amssymb}
\setcounter{secnumdepth}{-\maxdimen} % remove section numbering
\usepackage{iftex}
\ifPDFTeX
  \usepackage[T1]{fontenc}
  \usepackage[utf8]{inputenc}
  \usepackage{textcomp} % provide euro and other symbols
\else % if luatex or xetex
  \usepackage{unicode-math} % this also loads fontspec
  \defaultfontfeatures{Scale=MatchLowercase}
  \defaultfontfeatures[\rmfamily]{Ligatures=TeX,Scale=1}
\fi
\usepackage{lmodern}
\ifPDFTeX\else
  % xetex/luatex font selection
\fi
% Use upquote if available, for straight quotes in verbatim environments
\IfFileExists{upquote.sty}{\usepackage{upquote}}{}
\IfFileExists{microtype.sty}{% use microtype if available
  \usepackage[]{microtype}
  \UseMicrotypeSet[protrusion]{basicmath} % disable protrusion for tt fonts
}{}
\makeatletter
\@ifundefined{KOMAClassName}{% if non-KOMA class
  \IfFileExists{parskip.sty}{%
    \usepackage{parskip}
  }{% else
    \setlength{\parindent}{0pt}
    \setlength{\parskip}{6pt plus 2pt minus 1pt}}
}{% if KOMA class
  \KOMAoptions{parskip=half}}
\makeatother
\setlength{\emergencystretch}{3em} % prevent overfull lines
\providecommand{\tightlist}{%
  \setlength{\itemsep}{0pt}\setlength{\parskip}{0pt}}
\usepackage{tikz}
\usetikzlibrary{arrows.meta,calc,positioning}
\usepackage{bookmark}
\IfFileExists{xurl.sty}{\usepackage{xurl}}{} % add URL line breaks if available
\urlstyle{same}
\hypersetup{
  pdftitle={Classical Control Theory notes},
  pdfauthor={Samo F.},
  hidelinks,
  pdfcreator={LaTeX via pandoc}}

\title{Classical Control Theory notes}
\author{Samo F.}
\date{2025}

\begin{document}
\maketitle

\section{Introduction}\label{introduction}

Control system is mechanism that alters the future state of a system.
Control theory is a strategy to select the approriate inputs.

You could just try out all the values, but that's not a particulary
efficient method.

Control theory is converting unstable open loop systems to stable closed
loop systems, which are really just open loop systems again :). Common
definition of control theory includes expected/target value of the
output, but that's not really the case - control theory is modifying the
behaviour of the system so that expected value is achievable under given
constraints.

Inputs to the process of designing a control system are:

\begin{itemize}
\tightlist
\item
  objective
\item
  criteria (goals)
\item
  information about the system
\item
  current state of the system
\item
  constraints
\end{itemize}

\section{Systems and signals}\label{systems-and-signals}

Processes:

\begin{itemize}
\tightlist
\item
  continuous
\item
  discrete
\item
  batch
\end{itemize}

Energy of signal:

\[ E = \int\limits\_{t_1}^{t_2} |x(t)|^2\, dt \]

(Average) power of signal:

\[ P = \frac{1}{t*2 - t_1} \int\limits*{t_1}^{t_2} |x(t)|^2\, dt \]

Shannon:

\[ f > 2f_m \]

Signals:

\begin{itemize}
\tightlist
\item
  constant
\item
  linearly increasing
\item
  step \(1(t) = u(t)\)
\item
  exponential
\item
  sinusoidal
\item
  exponentially damped sinusoidal
\item
  square impulse
\item
  Dirac delta \(\int_{-\infty}^{\infty} \delta(t)\, dt = 1\)
\item
  impulse train
\end{itemize}

\section{Modeling}\label{modeling}

All models are wrong, but some are useful; most commonly for prediction,
control, optimization and training.

Best model should cover only the essential aspects of the system,
depending on the objective of the modelling.

Models usually start from LTI (linear time invariant) descriptions and
then are extended to include nonlinearities, time variance,
stochasticity, etc. And event then we try to get back to LTI usually by
linearization around the operating point (small signal response).

Conditions for LTI are:

\begin{itemize}
\tightlist
\item
  superposition

  \begin{itemize}
  \tightlist
  \item
    additivity \(T(x_1 + x_2) = T(x_1) + T(x_2)\)
  \item
    homogeneity \(T(ax) = aT(x)\)
  \end{itemize}
\item
  time invariance
\end{itemize}

Model can be made from purely theoritical first principles, from data
(system identification) or from a combination of the two.

.. see \href{AS.md}{AS notes}.. see book pp.~86 (Primer 3.1 Modeliranje
avtomobilskega vzmetenja)

\subsection{Laplace transform}\label{laplace-transform}

LTI systems can be described by linear differential equations, which can
be solved either painfully or with Laplace transform.

\[ \mathcal{L}\{f(t)\} = F(s) = \int\limits\_{0}^{\infty} e^{-st} f(t)\, dt \]

\section{Sources}\label{sources}

https://www.youtube.com/playlist?list=PLUMWjy5jgHK1NC52DXXrriwihVrYZKqjk

\end{document}
